\bookchapter{Глава 41. Панорама ИИ-инструментов для программирования в 2026 году}

\chapterintro{Сразу предупредим: эта глава — «снимок» экосистемы инструментов на август 2026 года. Через полгода он устареет — и это первый урок, который хочет преподать глава. Её цель не в том, чтобы научить «каким инструментом пользоваться», а в том, чтобы помочь выстроить рамку суждения «как выбирать инструменты». Понимая философию дизайна разных инструментов, вы в любой момент сделаете выбор, подходящий именно вам.}

\booksection{Три школы и рамка выбора}

ИИ-инструменты для программирования в 2026 году делятся примерно на три школы, за каждой стоит своя философия дизайна. Сначала взгляните на эту панораму:

\begin{center}
\begin{tabularx}{\linewidth}{|X|X|X|X|}
\hline
\textbf{Школа} & \textbf{Представители} & \textbf{Философия дизайна} & \textbf{Подходящие сценарии} \tabularnewline
\hline
ИИ, встроенный в IDE & Cursor 3, Trae, Windsurf, CodeBuddy & Глубокая интеграция ИИ-возможностей в редактор & Ежедневная разработка, ревью кода, постепенное улучшение \tabularnewline
Терминальный Agent & Claude Code, Codex CLI, Aider & Терминал как интерфейс: ИИ напрямую оперирует файловой системой и командами & Операции уровня проекта, массовые правки, интеграция в CI/CD \tabularnewline
Облачный Builder & v0, Bolt, Replit Agent & Создание целого проекта с нуля в браузере & Быстрые прототипы, старт без настроек, мгновенный предпросмотр \tabularnewline
\hline
\end{tabularx}
\end{center}

\begin{center}\small\color{gray}Три школы ИИ-инструментов программирования\end{center}

Ключ к выбору инструмента — не «какой лучший», а «какой лучше подходит вашему текущему сценарию». Три практичных критерия:

\begin{itemize}
  \item Каков ваш уровень в программировании? Начинаете с нуля — берите облачный Builder, есть база — ИИ в IDE, продвинутый уровень — терминальный Agent.
  \item На каком этапе ваш проект? Начинаете с нуля — облачный Builder, итеративная разработка — ИИ в IDE, поддержка и рефакторинг — терминальный Agent.
  \item Каковы ваши потребности в деплое? Быстрая проверка идей — облачный Builder, серьёзный проект — ИИ в IDE, автоматизация процессов — терминальный Agent.
\end{itemize}

\booksection{Школа ИИ, встроенного в IDE}

ИИ, встроенный в IDE, — самая массовая школа: ИИ-возможности встраиваются в ваш редактор кода, добавляя ИИ-усиление поверх классического опыта программирования. В 2026 году представители этой школы — Cursor 3, Trae, Windsurf и CodeBuddy.

\booksubsection{Cursor 3}

Cursor — сегодня самая широко используемая в мире ИИ-IDE. Вышедший в апреле 2026 года Cursor 3 стал важным поворотным пунктом для ИИ-инструментов: доля режима Agent впервые превысила режим автодополнения Tab — разработчики перешли от «ИИ дополняет мне код» к «ИИ выполняет за меня задачи». Ключевые особенности Cursor 3: редактирование нескольких файлов (ИИ может одновременно менять несколько файлов, выполняя задачу), понимание контекста (автоматически разбирается в структуре проекта и зависимостях), глубокая интеграция с терминалом (ИИ может запускать команды и смотреть вывод прямо в редакторе).

\booksubsection{Trae}

Trae — ИИ-IDE от ByteDance с большой пользовательской базой на китайском рынке. Его режим SOLO (он же режим Builder) — типичный пример возможности «создай проект одной фразой»: вы описываете требования, а ИИ сам проходит весь путь от инициализации проекта до генерации кода. Сильные стороны Trae — поддержка китайского языка и адаптация к сетевой среде Китая; подходит для быстрого старта пользователей внутри страны.

\booksubsection{Другие инструменты}

Windsurf (бывший Codeium) славится плавным автодополнением кода и подходит разработчикам, которые любят печатать «на лету». Tencent Cloud CodeBuddy, плагин для VS Code, силён в интеграции с облачной экосистемой — для проектов, планирующих деплой на Tencent Cloud или разработку мини-программ WeChat. Alibaba Tongyi Lingma и Baidu Comate предлагают похожий опыт интеграции с IDE, каждый со своими особенностями в оптимизации под китайскоязычные сценарии.

\booksection{Школа терминальных агентов}

Терминальные агенты — самая быстрорастущая школа 2026 года. Ключевая идея: ИИ не должен быть просто плагином редактора — он должен работать со всей средой разработки как разработчик: читать и писать файлы, запускать команды, смотреть вывод, ставить зависимости. Представитель школы — Claude Code.

\booksubsection{Claude Code}

Claude Code — терминальный ИИ-инструмент для программирования от Anthropic, к середине 2026 года достигший годового дохода в 2,5 млрд долларов США (в годовом исчислении). Работает он совершенно не так, как традиционная IDE: вы описываете задачу на естественном языке в терминале, а Claude Code сам читает нужные файлы, правит код, запускает тесты, смотрит результат, итеративно исправляется — и всё это без того, чтобы вы открывали редактор.

Сила Claude Code — «видение всего проекта»: его не ограничивает окно редактора, он свободно просматривает весь каталог проекта и понимает зависимости между файлами. Он также поддерживает файл памяти проекта CLAUDE.md, чтобы ИИ с начала каждой сессии знал контекст. Для рефакторинга больших проектов, массовых правок, написания тестов и подобных задач терминальный агент эффективнее плагина IDE.

\booksubsection{Codex CLI и Aider}

Codex CLI от OpenAI предлагает похожий опыт программирования в терминале и поддерживает функцию /import — можно импортировать существующую кодовую базу, чтобы AI её изучил. Aider — инструмент с открытым исходным кодом для AI-программирования в терминале; он поддерживает разные модели (GPT, Claude, DeepSeek и др.) и подходит тем, кто хочет гибко переключаться между моделями. Хотя у этих инструментов свои особенности, главная идея у них одна: дать AI в терминале полный контроль над средой разработки.

\booksection{Облачные конструкторы (Builder)}

Облачные конструкторы — самый подходящий вариант для новичков с нуля. Вам не нужно устанавливать никаких программ: достаточно открыть браузер и описать свою задачу, а AI создаст полный проект в облаке и даст мгновенный предпросмотр.

\begin{itemize}
  \item v0 (от Vercel): отлично генерирует интерфейсные UI-компоненты и страницы, подходит для быстрого создания прототипов веб-страниц.
  \item Bolt: позволяет создать полноценное full-stack-приложение из одной фразы — с интерфейсом, серверной частью и базой данных.
  \item Replit Agent: работает в облачной IDE Replit — не только создаёт проект, но и запускает и развёртывает его прямо в облаке.
\end{itemize}

Общее преимущество этих инструментов — «старт без настройки»: достаточно открыть браузер. Ограничение в том, что когда проект усложняется, облачной среде редактирования не хватает гибкости локальной IDE. Поэтому они лучше всего подходят для этапа прототипа «от 0 до 1»: быстро проверьте идею с помощью облачного конструктора, а после успешной проверки перенесите проект в локальную IDE для глубокой разработки.

\booksection{Инструменты меняются — методология остаётся}

Если завтра исчезнут все инструменты ИИ-программирования, должны остаться три привычки: сначала подумай, потом делай; за один раз — один маленький шаг; сделав — проверь. Инструменты — усилители, методология — фундамент. Усилители можно заменить, фундамент шататься не должен.

Вспомним, как менялись инструменты с 2024 по 2026 год: в 2024-м королём был GitHub Copilot, в 2025-м его сменил Cursor, а в 2026-м стремительно ворвался Claude Code. С появлением каждого нового инструмента кто-нибудь заявлял: «В этот раз всё по-другому», — но в итоге выяснялось: хорошо освоить новый инструмент удавалось именно тем, у кого прочная методология. Им на смену инструмента хватало нескольких дней адаптации, а тем, кто выучил только кнопки, приходилось начинать всё заново.

Поэтому рекомендуемый в этой книге путь обучения такой: начните с одного инструмента (подойдёт любой) и на практике прочувствуйте пять ключевых принципов мышления; когда методология станет достаточно крепкой, попробуйте инструменты разных школ и почувствуйте разницу в их философии; в конце концов у вас сложится свой набор инструментов — например, Cursor для повседневной разработки, Claude Code для рефакторинга и тестирования, v0 для быстрых прототипов. Набор у каждого свой, но методология универсальна.

\begin{tcbpractice}{Попробуйте сами: соберите себе «боевой набор» инструментов}
Определите своё место по трём критериям (уровень программирования, этап проекта, требования к развёртыванию): если вы новичок с нуля попробуйте Bolt или v0 и сделайте страницу; есть основа выполните одну итерационную задачу в Cursor или Trae; хотите вызов запустите Claude Code в терминале. А затем ответьте: какая школа лучше всего соответствует вашему текущему этапу?

\end{tcbpractice}

\begin{tcbsummary}{Итоги главы}
\begin{itemize}
  \item Три школы: встроенный в IDE AI (повседневная разработка), терминальный Agent (операции на уровне проекта), облачный Builder (прототипы с нуля)
  \item Выбирая инструмент, задайте три вопроса: каков уровень программирования? на каком этапе проект? каковы требования к развёртыванию?
  \item В Cursor 3 режим Agent впервые обогнал Tab-дополнение; годовая выручка Claude Code достигла 2,5 млрд долларов
  \item Облачный Builder подходит для этапа «от 0 до 1»; после проверки переносите проект в локальную IDE для глубокой разработки
  \item Инструменты обновляются каждые три месяца: тот, кто освоил методологию, адаптируется к новому инструменту за несколько дней
\end{itemize}

\end{tcbsummary}
